\documentclass[aspectratio=169, 11pt]{beamer}

% ── Tema i colors ──────────────────────────────────────────
\usetheme{Madrid}
\usecolortheme{default}

% Colors personalitzats — to càlid, acadèmic
\definecolor{uocblau}{RGB}{0, 74, 124}
\definecolor{uocgris}{RGB}{100, 100, 100}
\definecolor{fonsclar}{RGB}{250, 250, 250}

\setbeamercolor{palette primary}{bg=uocblau, fg=white}
\setbeamercolor{palette secondary}{bg=uocblau!80, fg=white}
\setbeamercolor{palette tertiary}{bg=uocblau!60, fg=white}
\setbeamercolor{frametitle}{bg=uocblau, fg=white}
\setbeamercolor{title}{fg=uocblau}
\setbeamercolor{item}{fg=uocblau}
\setbeamercolor{block title}{bg=uocblau, fg=white}

% Treure elements de navegació innecessaris
\setbeamertemplate{navigation symbols}{}
\setbeamertemplate{footline}[frame number]

% ── Paquets ────────────────────────────────────────────────
\usepackage[utf8]{inputenc}
\usepackage[T1]{fontenc}
\usepackage[catalan]{babel}
\usepackage{graphicx}
\usepackage{hyperref}
\usepackage{booktabs}
\usepackage{multicol}

% ── Metadades ──────────────────────────────────────────────
\title{PAC 2 --- Visualització de Dades}
\subtitle{Pie Chart · Nightingale Rose Chart · Bump Chart}
\author{Ivan Rodríguez Quintana}
\institute{Universitat Oberta de Catalunya}
\date{Març 2026}

% ══════════════════════════════════════════════════════════════
\begin{document}

% ── PORTADA ────────────────────────────────────────────────
\begin{frame}[plain]
\centering
\vspace{1.5em}

{\color{uocblau}\rule{\textwidth}{2pt}}

\vspace{1em}
{\Large\bfseries\color{uocblau} PAC 2 --- Visualització de Dades}

\vspace{0.5em}
{\large Pie Chart · Nightingale Rose Chart · Bump Chart}

\vspace{1em}
{\color{uocblau}\rule{\textwidth}{0.5pt}}

\vspace{1.5em}
{\normalsize Ivan Rodríguez Quintana}\\[0.3em]
{\small Universitat Oberta de Catalunya}\\[0.3em]
{\small Març 2026}

\vspace{1.5em}
{\color{uocblau}\rule{\textwidth}{2pt}}
\end{frame}

% ── INTRODUCCIÓ ────────────────────────────────────────────
\begin{frame}{Objectiu de la PAC}
\begin{itemize}
    \item Crear \textbf{3 visualitzacions interactives} amb tècniques assignades
    \item Publicar-les amb \textbf{URL pública} accessible sense registre
    \item Demostrar comprensió de cada tècnica i coherència dades--visualització
\end{itemize}

\vspace{1em}

\begin{block}{Eines utilitzades}
    \textbf{Python 3} + \textbf{Plotly} (graph\_objects i express) + \textbf{Pandas}\\
    Publicació: \textbf{GitHub Pages} \hfill
    \footnotesize{\url{https://rebanyat.github.io/VD-PAC2/}}
\end{block}
\end{frame}

% ── TÈCNIQUES ASSIGNADES ──────────────────────────────────
\begin{frame}{Tècniques i datasets}
\begin{table}
\centering
\small
\begin{tabular}{@{}lll@{}}
\toprule
\textbf{Tècnica} & \textbf{Dataset} & \textbf{Font} \\
\midrule
Pie Chart & Formes d'OVNIs (78.400 reg.) & NUFORC / Kaggle \\
Nightingale Rose & Impactes aus--avions (25.429 reg.) & FAA / Kaggle \\
Bump Chart & Consum xocolata per càpita & FAO / Our World in Data \\
\bottomrule
\end{tabular}
\end{table}
\end{frame}

% ══════════════════════════════════════════════════════════════
% SECCIÓ 1: PIE CHART
% ══════════════════════════════════════════════════════════════

\begin{frame}{Pie Chart --- Definició i funcionament}
\begin{columns}
\begin{column}{0.55\textwidth}
\begin{itemize}
    \item Gràfic circular que divideix un \textbf{total} en sectors proporcionals
    \item Cada sector representa una \textbf{categoria} i la seva proporció sobre el total
    \item L'angle del sector és proporcional al valor que representa
    \item Adequat per a \textbf{poques categories} (5--10 màx.)
\end{itemize}
\end{column}
\begin{column}{0.4\textwidth}
\begin{block}{Limitacions}
    \small
    \begin{itemize}
        \item Difícil comparar sectors de mida similar
        \item Perd llegibilitat amb moltes categories
        \item L'ull humà compara malament àrees/angles
    \end{itemize}
\end{block}
\end{column}
\end{columns}
\end{frame}

\begin{frame}{Pie Chart --- Dataset i transformació}
\begin{block}{Dataset: NUFORC (National UFO Reporting Center)}
    78.400 avistaments d'OVNIs als EUA, registrats entre 1949 i 2014.
\end{block}

\vspace{0.8em}

\textbf{Transformació aplicada:}
\begin{enumerate}
    \item Eliminar registres sense forma reportada (1.932 files)
    \item Agrupar per la columna \texttt{shape} i comptar freqüències
    \item Mantenir les \textbf{10 formes més freqüents}, resta com a ``Diverses formes''
    \item Traducció dels noms de formes de l'anglès al català
\end{enumerate}

\vspace{0.5em}
\small
\textbf{Resultat:} 11 categories (10 + Diverses) --- ideal per a un pie chart.
\end{frame}

\begin{frame}{Pie Chart --- Visualització i anàlisi}
\centering
\textit{[Demostració en viu: \url{https://rebanyat.github.io/VD-PAC2/pie_chart.html}]}

\vspace{1.5em}

\begin{columns}
\begin{column}{0.5\textwidth}
\textbf{Observacions clau:}
\begin{itemize}
    \item ``Llum'' domina amb un 21\% dels avistaments
    \item Les 3 formes principals (Llum, Triangle, Cercle) sumen quasi el 41\%
    \item ``Bola de foc'' és la 4a forma més comuna
\end{itemize}
\end{column}
\begin{column}{0.45\textwidth}
\textbf{Per què encaixa:}
\begin{itemize}
    \item Poques categories ben diferenciades
    \item Objectiu: mostrar proporcions d'un total
    \item La temàtica còsmica reforça el contingut visual
\end{itemize}
\end{column}
\end{columns}
\end{frame}

% ══════════════════════════════════════════════════════════════
% SECCIÓ 2: NIGHTINGALE ROSE CHART
% ══════════════════════════════════════════════════════════════

\begin{frame}{Nightingale Rose Chart --- Origen i funcionament}
\begin{columns}
\begin{column}{0.55\textwidth}
\begin{itemize}
    \item Inventada per \textbf{Florence Nightingale} el \textbf{1858}
    \item Usada per visualitzar causes de mortalitat a la Guerra de Crimea
    \item Gràfic polar on l'\textbf{àrea} del segment (no el radi) representa el valor
    \item Tots els segments tenen el \textbf{mateix angle}; el radi varia
\end{itemize}
\end{column}
\begin{column}{0.4\textwidth}
\begin{block}{Diferència clau vs.\ bar chart}
    \small
    Com que àrea $= \pi r^2 \cdot \frac{\theta}{2\pi}$ i $\theta$ és constant, cal aplicar $r = \sqrt{n}$ perquè l'àrea sigui proporcional al valor real.
    Sense la correcció, les diferències s'exageren visualment.
\end{block}
\end{column}
\end{columns}
\end{frame}

\begin{frame}{Nightingale Rose Chart --- Dataset i transformació}
\begin{block}{Dataset: FAA Wildlife Strike Database}
    25.429 incidents de col·lisió entre fauna i aeronaus als EUA (2000--2011).
\end{block}

\vspace{0.8em}

\textbf{Transformació aplicada:}
\begin{enumerate}
    \item Parsejar la columna \texttt{FlightDate} (format M/D/YY)
    \item Extreure el mes de cada incident
    \item Comptar incidents per mes, amb desglossament per danys
    \item Aplicar transformació $r = \sqrt{n}$ per a representació fidel de l'àrea
\end{enumerate}

\vspace{0.5em}
\small
\textbf{Resultat:} 12 valors mensuals amb 2 capes (amb/sense danys) --- perfecte per a un gràfic cíclic.
\end{frame}

\begin{frame}{Nightingale Rose Chart --- Visualització i anàlisi}
\centering
\textit{[Demostració en viu: \url{https://rebanyat.github.io/VD-PAC2/nightingale.html}]}

\vspace{1.5em}

\begin{columns}
\begin{column}{0.5\textwidth}
\textbf{Observacions clau:}
\begin{itemize}
    \item Pic clar a l'\textbf{Agost} (3.710 incidents)
    \item Mínim al \textbf{Febrer} (772 incidents)
    \item Patró estacional: coincideix amb la migració d'aus
\end{itemize}
\end{column}
\begin{column}{0.45\textwidth}
\textbf{Per què encaixa:}
\begin{itemize}
    \item Dades cícliques (12 mesos)
    \item Patró estacional clar i repetitiu
    \item Homenatge a l'ús original de Nightingale (mortalitat per mes)
\end{itemize}
\end{column}
\end{columns}
\end{frame}

% ══════════════════════════════════════════════════════════════
% SECCIÓ 3: BUMP CHART
% ══════════════════════════════════════════════════════════════

\begin{frame}{Bump Chart --- Definició i funcionament}
\begin{columns}
\begin{column}{0.55\textwidth}
\begin{itemize}
    \item Mostra com canvien els \textbf{rànkings} de diverses entitats al llarg del temps
    \item Eix X: temps --- Eix Y: posició al rànking (invertit, \#1 a dalt)
    \item Cada entitat és una línia connectada que ``salta'' amunt o avall
    \item \textbf{Color} per distingir cada entitat
\end{itemize}
\end{column}
\begin{column}{0.4\textwidth}
\begin{block}{Avantatges}
    \small
    \begin{itemize}
        \item Destaca canvis significatius de posició
        \item Detecta tendències (ascens, descens, estabilitat)
        \item Línies rectes = rànking estable
    \end{itemize}
\end{block}
\end{column}
\end{columns}
\end{frame}

\begin{frame}{Bump Chart --- Dataset i transformació}
\begin{block}{Dataset: FAO via Our World in Data}
    Consum de cacau per càpita (kg/any) per país, de 1961 a 2023.
\end{block}

\vspace{0.8em}

\textbf{Transformació aplicada:}
\begin{enumerate}
    \item Eliminar regions agregades (FAO, continents) i nacions petites
    \item Filtrar període 1990--2023
    \item Seleccionar els 10 països amb major consum mitjà
    \item Calcular rànking per any ($1 =$ consum més alt)
    \item Traducció dels noms de països al català
\end{enumerate}

\vspace{0.5em}
\small
\textbf{Resultat:} 10 països × 34 anys --- ideal per a un bump chart dinàmic.
\end{frame}

\begin{frame}{Bump Chart --- Visualització i anàlisi}
\centering
\textit{[Demostració en viu: \url{https://rebanyat.github.io/VD-PAC2/bump_chart.html}]}

\vspace{1.5em}

\begin{columns}
\begin{column}{0.5\textwidth}
\textbf{Observacions clau:}
\begin{itemize}
    \item \textbf{Costa d'Ivori} lidera el 2023 --- sorprenent, ja que és productor
    \item \textbf{Luxemburg} i \textbf{Islàndia} es mantenen estables al top
    \item Canvis significatius al llarg de les dècades
\end{itemize}
\end{column}
\begin{column}{0.45\textwidth}
\textbf{Per què encaixa:}
\begin{itemize}
    \item Dades de rànking que evolucionen en el temps
    \item 10 entitats: suficients per veure competència
    \item Interactivitat (hover, clic) ajuda la comprensió
\end{itemize}
\end{column}
\end{columns}
\end{frame}

% ══════════════════════════════════════════════════════════════
% TANCAMENT
% ══════════════════════════════════════════════════════════════

\begin{frame}{URLs públiques}
\centering
\vspace{1em}

\large
\begin{tabular}{@{}ll@{}}
\textbf{Índex:} & \url{https://rebanyat.github.io/VD-PAC2/} \\[0.8em]
\textbf{Pie Chart:} & \url{https://rebanyat.github.io/VD-PAC2/pie_chart.html} \\[0.8em]
\textbf{Nightingale:} & \url{https://rebanyat.github.io/VD-PAC2/nightingale.html} \\[0.8em]
\textbf{Bump Chart:} & \url{https://rebanyat.github.io/VD-PAC2/bump_chart.html} \\
\end{tabular}

\vspace{1.5em}
\normalsize
Totes les visualitzacions són \textbf{interactives} i accessibles sense registre.
\end{frame}

\begin{frame}{Conclusions}
\begin{itemize}
    \item Cada tècnica s'adapta a un \textbf{tipus de dades} específic:
    \begin{itemize}
        \item Pie chart $\rightarrow$ proporcions d'un total (poques categories)
        \item Nightingale $\rightarrow$ dades cícliques amb magnitud variable
        \item Bump chart $\rightarrow$ rànkings que evolucionen en el temps
    \end{itemize}

    \vspace{0.8em}
    \item La \textbf{transformació de dades} és clau: la correcció $r = \sqrt{n}$ al Nightingale evita exagerar diferències

    \vspace{0.8em}
    \item La \textbf{interactivitat} (hover, clic a la llegenda) millora significativament la comprensió de gràfics complexos com el bump chart

    \vspace{0.8em}
    \item L'elecció de \textbf{colors temàtics} (còsmic, estacional, xocolata) reforça el missatge de cada visualització
\end{itemize}
\end{frame}

\begin{frame}{}
\centering
\vspace{3em}
{\Large \textbf{Gràcies}}

\vspace{2em}
Ivan Rodríguez Quintana\\
\small Visualització de Dades --- UOC --- Març 2026

\vspace{1em}
\footnotesize
\url{https://rebanyat.github.io/VD-PAC2/}

\vspace{1em}
{\scriptsize\color{uocgris} Esbós inicial de les visualitzacions amb l'ajuda de Claude (Anthropic)}
\end{frame}

\end{document}
